\documentclass[12pt]{article}
\usepackage[margin=1in]{geometry}
\usepackage{amsmath, amssymb}
\usepackage{booktabs}
\usepackage{hyperref}
\usepackage{natbib}
\usepackage{graphicx}

\title{Signal Half-Life in Crypto Microstructure:\\
       An Empirical Study of Alpha Decay\\
       in Binance Perpetual Futures}

\author{Your Name\\
\small Independent Research\\
\small \href{https://github.com/YOUR_USERNAME/alpha-decay-lab}{github.com/YOUR\_USERNAME/alpha-decay-lab}}

\date{\today}

\begin{document}
\maketitle

\begin{abstract}
% TODO (Phase 5): Write abstract last, after results are final.
% Target: 150 words. Cover: motivation, data, methods, main findings, implications.
We study the rate at which microstructure trading signals lose predictive power
in cryptocurrency perpetual futures markets. Using tick-level data from Binance,
we estimate the half-life of three canonical signals --- order flow imbalance,
trade sign autocorrelation, and short-term price reversal --- across
Bitcoin and Ethereum perpetuals over [PERIOD]. 
We find that [MAIN RESULT]. We propose a Kalman-filter-based adaptive weighting
scheme that [PERFORMANCE CLAIM]. 
All code and data are available at the GitHub repository above.
\end{abstract}

\section{Introduction}
\label{sec:intro}

% TODO (Phase 5): ~1.5 pages
% Cover: economic motivation, why decay matters for practitioners,
% gap in the literature (most decay studies use equity data, not crypto),
% contribution of this paper.

Alpha decay is a central concern of quantitative trading. When a new
signal achieves widespread adoption, arbitrage activity erodes its
information content until the signal ceases to generate excess returns.
The rate of this erosion --- the signal's \emph{half-life} --- determines
how quickly strategies must adapt and how frequently signals must be
refreshed.

\section{Data}
\label{sec:data}

% TODO: fill in after running notebook 01 and 02.
% Cover: data source, time period, preprocessing steps, summary statistics.

We use publicly available trade data from Binance's REST API for
BTCUSDT and ETHUSDT perpetual futures contracts. The sample covers
[START DATE] to [END DATE], yielding approximately [N] million trades
per instrument. We discard observations within [X] minutes of daily
settlement to avoid settlement-driven distortions.

\section{Signal construction}
\label{sec:signals}

% TODO: fill in after running notebook 02.

We study three signals:

\subsection{Order flow imbalance}
% Reference Cont, Kukanov & Stoikov (2014)

\subsection{Trade sign autocorrelation}

\subsection{Short-term reversal}

\section{Methodology}
\label{sec:method}

\subsection{Information coefficient}

The Information Coefficient (IC) at horizon $h$ is the Spearman rank
correlation between the signal at time $t$ and the forward log return
$r_{t,t+h}$:

\begin{equation}
    \text{IC}_t = \rho_S\!\left(f_t,\, r_{t,t+h}\right)
\end{equation}

We compute a rolling IC with window $W$ to track decay over time.

\subsection{Half-life estimation}

We model IC decay as an exponential process:

\begin{equation}
    \text{IC}(t) = A \cdot e^{-\lambda t}
\end{equation}

The half-life is $\tau_{1/2} = \ln(2) / \lambda$. We estimate $\lambda$
via non-linear least squares \citep{scipy2020}.

\subsection{Adaptive weighting}

% TODO: describe Kalman filter approach from halflife/adaptive.py

\section{Results}
\label{sec:results}

% TODO (Phase 3): fill in after running notebook 06.
% Include: Table 1 (half-life by signal and instrument),
%          Figure 1 (IC decay curves),
%          Figure 2 (adaptive weight over time)

\begin{table}[h]
\centering
\caption{Estimated signal half-lives (periods)}
\label{tab:halflife}
\begin{tabular}{lrrrr}
\toprule
Signal & Instrument & Half-life & IC mean & IC IR \\
\midrule
OFI           & BTCUSDT & -- & -- & -- \\
OFI           & ETHUSDT & -- & -- & -- \\
Trade signs   & BTCUSDT & -- & -- & -- \\
Trade signs   & ETHUSDT & -- & -- & -- \\
Reversal      & BTCUSDT & -- & -- & -- \\
Reversal      & ETHUSDT & -- & -- & -- \\
\bottomrule
\end{tabular}
\end{table}

\section{Conclusion}
\label{sec:conclusion}

% TODO: ~0.5 pages.

\bibliographystyle{plainnat}
\bibliography{references}

\end{document}
